% ============================================================
\section{Observable-Relative Closure and Minimal Returned Data}
\label{sec:observable-relative-closure}
% ============================================================

\begin{quote}\small
\textbf{What this adds.}
The split and Schur calculus already determines the exact visible law and the
exact returned residual. What remains application-dependent is \emph{which}
visible outputs matter. This section makes that choice native: it records the
typed defect seen by a selected workload, the exact returned datum needed to
reconstruct that workload, and the finite-dimensional remainder seen after
Schur-jet truncation.

\smallskip\noindent
\textbf{Boundary.}
This is still operator-side calculus. It does not add new dynamics. It
composes the existing split and Schur operators with an explicit output map and
then studies the exact factorization profile of the resulting transfer.
\end{quote}

\subsection{Typed defects}

\begin{definition}[Typed coherence defect]
\label{def:typed-coherence-defect}\lean{CoherenceCalculus.typedCoherenceDefect}
Let \(X\) and \(Y\) be vector spaces, let \(P_X:X\to X\) and \(P_Y:Y\to Y\) be
projections, and let \(F:X\to Y\) be a map. The \emph{typed coherence defect}
of \(F\) at \(x\in X\) is
\[
\tau_{P_X,P_Y}(F;x):=P_YF(x)-F(P_Xx).
\]
\end{definition}

\begin{proposition}[Typed chain rule]
\label{prop:typed-chain-rule}\lean{CoherenceCalculus.typed_chain_rule}
Let \(F:Y\to Z\) and \(G:X\to Y\) be maps, and let \(P_X,P_Y,P_Z\) be
projections on \(X,Y,Z\), respectively. Then
\[
\tau_{P_X,P_Z}(F\circ G;x)
=
\tau_{P_Y,P_Z}(F;G(x))
+
\bigl(F(P_YG(x))-F(G(P_Xx))\bigr).
\]
If \(F\) is linear, the second term becomes \(F(\tau_{P_X,P_Y}(G;x))\).
\end{proposition}

\begin{proof}
\begin{align*}
\tau_{P_X,P_Z}(F\circ G;x)
&=
P_ZF(G(x))-F(G(P_Xx)) \\
&=
\bigl(P_ZF(G(x))-F(P_YG(x))\bigr)
+
\bigl(F(P_YG(x))-F(G(P_Xx))\bigr) \\
&=
\tau_{P_Y,P_Z}(F;G(x))
+
\bigl(F(P_YG(x))-F(G(P_Xx))\bigr).
\end{align*}
If \(F\) is linear, then
\[
F(P_YG(x))-F(G(P_Xx))
=
F\bigl(P_YG(x)-G(P_Xx)\bigr)
=
F(\tau_{P_X,P_Y}(G;x)).
\]
\end{proof}

\begin{definition}[Typed bilinear defect]
\label{def:typed-bilinear-defect}\lean{CoherenceCalculus.typedBilinearDefect}
Let \(X,Y,Z\) be vector spaces, let \(P_X,P_Y,P_Z\) be projections, and let
\(B:Y\times X\to Z\) be bilinear. The \emph{typed bilinear defect} is
\[
\tau_{P_Y,P_X,P_Z}(B;y,x):=
P_ZB(y,x)-B(P_Yy,P_Xx).
\]
\end{definition}

\begin{proposition}[Typed bilinear defect decomposition]
\label{prop:typed-bilinear-defect-decomposition}\lean{CoherenceCalculus.typed_bilinear_defect_decomposition}
Let \(Q_X:=I-P_X\) and \(Q_Y:=I-P_Y\). Assume the visible-visible closure
condition
\[
P_ZB(P_Yy,P_Xx)=B(P_Yy,P_Xx)
\qquad (x\in X,\ y\in Y).
\]
Then
\[
\tau_{P_Y,P_X,P_Z}(B;y,x)
=
P_ZB(P_Yy,Q_Xx)
+
P_ZB(Q_Yy,P_Xx)
+
P_ZB(Q_Yy,Q_Xx).
\]
\end{proposition}

(Proof in Appendix~\ref{sec:appendix-elimination}.)

\subsection{Observable packages and transfer operators}

\begin{definition}[Observable package]
\label{def:observable-package}\lean{CoherenceCalculus.ObservablePackage}
Let \(X\) be a vector space. An \emph{observable package} on \(X\) is a triple
\[
\Omega=(S,V,O)
\]
consisting of a source subspace \(S\subseteq X\), an output space \(V\), and a
linear observable map \(O:X\to V\).
\end{definition}

\begin{definition}[Paired observable package]
\label{def:paired-observable-package}\lean{CoherenceCalculus.PairedObservablePackage}
Let \(X\) and \(Y\) be vector spaces. A \emph{paired observable package} is a
triple
\[
\Omega=(S,T,B)
\]
consisting of a source subspace \(S\subseteq X\), a test subspace
\(T\subseteq Y\), and a bilinear pairing
\[
B:Y\times X\to \mathbb{R}
\qquad\text{or}\qquad
B:Y\times X\to \mathbb{C}.
\]
It induces the linear observable map
\[
O_T:X\to T^\ast,
\qquad
O_T(x):=B(\,\cdot\,,x)\big|_T.
\]
All linear observable-transfer statements below apply to \(O_T\).
\end{definition}

\begin{definition}[Direct observable transfer]
\label{def:direct-observable-transfer}\lean{CoherenceCalculus.directObservableTransfer}
Let \(\Omega=(S,V,O)\) be an observable package on \(X\), let \(P\) be a
projection on \(X\) with complement \(Q:=I-P\), and let \(A:X\to X\) be
linear. The \emph{direct observable transfer} is the map
\[
\mathfrak D_{\Omega}(P,A):=
OAQ\big|_{S}:S\to V.
\]
\end{definition}

\begin{proposition}[Exact direct observable decomposition]
\label{prop:direct-observable-decomposition}\lean{CoherenceCalculus.direct_observable_decomposition}
In the setting of Definition~\ref{def:direct-observable-transfer},
\[
O(Ax)=O(APx)+\mathfrak D_{\Omega}(P,A)x
\qquad (x\in S).
\]
In particular, if \(x\in S\cap PX\), then \(\mathfrak D_{\Omega}(P,A)x=0\).
\end{proposition}

\begin{proof}
Because \(I=P+Q\),
\[
Ax=A(P+Q)x=APx+AQx.
\]
Applying \(O\) gives
\[
O(Ax)=O(APx)+OAQx=O(APx)+\mathfrak D_{\Omega}(P,A)x.
\]
If \(x\in PX\), then \(Qx=0\), so the last term vanishes.
\end{proof}

\begin{definition}[Observable Schur return transfer]
\label{def:observable-schur-return-transfer}\lean{CoherenceCalculus.observableSchurReturnTransfer}
Let \(\Omega=(S,V,O)\) be an observable package on the Hilbert space
\(\cH\), let \(P\) be an orthogonal projection, and assume
\(\CCSigmaP{P}{z}\) is defined. On the visible workload
\[
S_P:=S\cap P\cH
\]
define the \emph{observable Schur return transfer} by
\[
\mathfrak R_{\Omega}(P,A;z)
:=
O\,\CCSigmaP{P}{z}\big|_{S_P}:S_P\to V.
\]
\end{definition}

\begin{definition}[Fixed-source observable Schur return transfer]
\label{def:fixed-source-observable-schur-return-transfer}\lean{CoherenceCalculus.fixedSourceObservableSchurReturnTransfer}
Let \(\Omega=(S,V,O)\) be an observable package on the Hilbert space
\(\cH\), let \(P\) be an orthogonal projection, and assume
\(\CCSigmaP{P}{z}\) is defined. The \emph{fixed-source observable Schur return
transfer} is
\[
\widetilde{\mathfrak R}_{\Omega}(P,A;z)
:=
O\,\CCSigmaP{P}{z}\,P\big|_{S}:S\to V.
\]
\end{definition}

\begin{proposition}[Returned residual seen by the workload]
\label{prop:observable-return-is-schur}\lean{CoherenceCalculus.observable_return_is_schur}
In the setting of Definition~\ref{def:observable-schur-return-transfer}, every
\(u\in S_P\) satisfies
\[
O\bigl(\mathcal F^{\mathrm{return}}_{P,A,z}(u)\bigr)
=
\mathfrak R_{\Omega}(P,A;z)u,
\]
and
\[
O\bigl(\CCAeff{A}{P}(z)u\bigr)
=
O(PAPu)+z\,\mathfrak R_{\Omega}(P,A;z)u.
\]
\end{proposition}

\begin{proof}
The first identity is Proposition~\ref{prop:return-residual-is-schur} composed
with \(O\):
\[
O\bigl(\mathcal F^{\mathrm{return}}_{P,A,z}(u)\bigr)
=
O\bigl(\CCSigmaP{P}{z}u\bigr)
=
\mathfrak R_{\Omega}(P,A;z)u.
\]
For the second identity, Definition~\ref{def:effective-operator} gives
\[
\CCAeff{A}{P}(z)u=PAPu+z\,\CCSigmaP{P}{z}u.
\]
Applying \(O\) and using the first identity yields the claim.
\end{proof}

\begin{proposition}[Fixed-source return restricts to the visible workload]
\label{prop:fixed-source-return-restricts}\lean{CoherenceCalculus.fixed_source_return_factorization, CoherenceCalculus.fixed_source_return_on_fixed_points}
In the setting of
Definitions~\ref{def:observable-schur-return-transfer}
and~\ref{def:fixed-source-observable-schur-return-transfer}, every
\(u\in S_P\) satisfies
\[
\widetilde{\mathfrak R}_{\Omega}(P,A;z)u
=
\mathfrak R_{\Omega}(P,A;z)u.
\]
If in addition \(P(S)\subseteq S\), then \(P\big|_S:S\to S_P\) and
\[
\widetilde{\mathfrak R}_{\Omega}(P,A;z)
=
\mathfrak R_{\Omega}(P,A;z)\circ P\big|_S.
\]
Consequently,
\[
\Img\bigl(\widetilde{\mathfrak R}_{\Omega}(P,A;z)\bigr)
=
\Img\bigl(\mathfrak R_{\Omega}(P,A;z)\bigr),
\]
hence the two return transfers have the same rank.
\end{proposition}

\begin{proof}
If \(u\in S_P\), then \(Pu=u\). Therefore
\[
\widetilde{\mathfrak R}_{\Omega}(P,A;z)u
=
O\,\CCSigmaP{P}{z}Pu
=
O\,\CCSigmaP{P}{z}u
=
\mathfrak R_{\Omega}(P,A;z)u.
\]
If \(P(S)\subseteq S\), then \(Px\in S_P\) for every \(x\in S\), so the
displayed factorization is exactly the identity just proved, applied to
\(u:=Px\).

Under the same hypothesis,
\(\Img(\widetilde{\mathfrak R}_{\Omega}(P,A;z))\subseteq
\Img(\mathfrak R_{\Omega}(P,A;z))\) by the factorization. For the reverse
inclusion, let \(y=\mathfrak R_{\Omega}(P,A;z)u\) with \(u\in S_P\). Since
\(u\in S\) and \(Pu=u\),
\[
y
=
\widetilde{\mathfrak R}_{\Omega}(P,A;z)u,
\]
so \(y\in \Img(\widetilde{\mathfrak R}_{\Omega}(P,A;z))\).
\end{proof}

\begin{remark}[Two return-transfer roles]
\label{rem:two-return-transfer-roles}\lean{CoherenceCalculus.observableSchurReturnTransfer, CoherenceCalculus.fixedSourceObservableSchurReturnTransfer, CoherenceCalculus.fixed_source_return_factorization, CoherenceCalculus.fixed_source_return_on_fixed_points}
The map \(\mathfrak R_{\Omega}(P,A;z)\) is the pointwise return seen on the
actual visible workload \(S_P=S\cap P\cH\). The map
\(\widetilde{\mathfrak R}_{\Omega}(P,A;z)\) keeps the source space fixed at
\(S\), so it is the correct object for comparing different cuts, taking ranks
and singular values across cuts, and formulating continuity on the
Grassmannian.
\end{remark}

\begin{proposition}[Continuity of the fixed-source return transfer]
\label{prop:fixed-source-return-continuity}\lean{CoherenceCalculus.FixedSourceReturnContinuityData, CoherenceCalculus.fixed_source_return_continuity}
Let \(\Omega=(S,V,O)\) be an observable package on a finite-dimensional
Hilbert space \(\cH\), let \(\mathcal U\) be a family of orthogonal
projections, and fix \(z\). Assume \(P\mapsto \CCSigmaP{P}{z}\) is continuous
on \(\mathcal U\) in operator norm. Then
\[
P\longmapsto \widetilde{\mathfrak R}_{\Omega}(P,A;z)
\]
is continuous on \(\mathcal U\) as a map into the operator space
\(\mathcal B(S,V)\).
\end{proposition}

\begin{proof}
Let \(P_n\to P\) in \(\mathcal U\). By hypothesis,
\(\CCSigmaP{P_n}{z}\to \CCSigmaP{P}{z}\) in operator norm. Therefore
\begin{align*}
\bigl\|
\widetilde{\mathfrak R}_{\Omega}(P_n,A;z)
-\widetilde{\mathfrak R}_{\Omega}(P,A;z)
\bigr\|_{\mathrm{op}}
&=
\bigl\|
O\bigl(\CCSigmaP{P_n}{z}P_n-\CCSigmaP{P}{z}P\bigr)\big|_S
\bigr\|_{\mathrm{op}} \\
&\le
\|O\|\,
\bigl\|
\CCSigmaP{P_n}{z}P_n-\CCSigmaP{P}{z}P
\bigr\|_{\mathrm{op}} \\
&\le
\|O\|\,
\Bigl(
\|\CCSigmaP{P_n}{z}-\CCSigmaP{P}{z}\|_{\mathrm{op}}\|P_n\|_{\mathrm{op}}
+ \|\CCSigmaP{P}{z}\|_{\mathrm{op}}\|P_n-P\|_{\mathrm{op}}
\Bigr) \\
&\le
\|O\|\,
\Bigl(
\|\CCSigmaP{P_n}{z}-\CCSigmaP{P}{z}\|_{\mathrm{op}}
+ \|\CCSigmaP{P}{z}\|_{\mathrm{op}}\|P_n-P\|_{\mathrm{op}}
\Bigr),
\end{align*}
because orthogonal projections have operator norm \(1\). The right-hand side
tends to \(0\), which proves continuity.
\end{proof}

\begin{definition}[Task-perfect coherence]
\label{def:task-perfect-coherence}\lean{CoherenceCalculus.taskPerfectCoherence}
Let \(\Omega=(S,V,O)\) be an observable package on \(\cH\). The pair
\((A,P)\) is \emph{task-perfect for \(\Omega\)} if
\[
\widetilde{\mathfrak R}_{\Omega}(P,A;z)=0
\]
for every \(z\) for which the Schur correction is defined.
\end{definition}

\begin{corollary}[Whole-state perfect coherence implies task-perfect coherence]
\label{cor:perfect-implies-task-perfect}\lean{CoherenceCalculus.perfect_implies_taskPerfect}
If \(\CCSigmaP{P}{z}=0\) for every \(z\) where defined, then \((A,P)\) is
task-perfect for every observable package \(\Omega=(S,V,O)\).
\end{corollary}

\begin{proof}
If \(\CCSigmaP{P}{z}=0\), then
\[
\widetilde{\mathfrak R}_{\Omega}(P,A;z)
=
O\,\CCSigmaP{P}{z}\,P\big|_{S}
=
0.
\]
\end{proof}

\subsection{Task-visible defect packages}

\begin{definition}[Source-visible loss]
\label{def:source-visible-loss}\lean{CoherenceCalculus.sourceVisibleLoss}
Let \(\Omega=(S,V,O)\) be an observable package on \(X\), and let \(P\) be a
projection on \(X\). The \emph{source-visible loss} of the cut \(P\) relative
to \(\Omega\) is
\[
\mathfrak S_{\Omega}(P):=O(I-P)\big|_S:S\to V.
\]
\end{definition}

\begin{definition}[Task-visible defect package]
\label{def:task-visible-defect-package}\lean{CoherenceCalculus.TaskVisibleDefectPackage, CoherenceCalculus.taskVisibleDefectPackage}
Let \(\Omega=(S,V,O)\) be an observable package on \(\cH\), let \(P\) be an
orthogonal projection, and assume \(\CCSigmaP{P}{z}\) is defined. The
\emph{task-visible defect package} is the triple
\[
\mathcal E_{\Omega}(P,A;z)
:=
\Bigl(
\mathfrak S_{\Omega}(P),\,
\mathfrak D_{\Omega}(P,A),\,
\widetilde{\mathfrak R}_{\Omega}(P,A;z)
\Bigr).
\]
\end{definition}

\begin{definition}[Task-visible faithfulness classes]
\label{def:task-visible-faithfulness}\lean{CoherenceCalculus.TaskVisibleFaithfulness}
Let \(\Omega=(S,V,O)\) be an observable package on \(\cH\).
\begin{enumerate}[label=(\roman*), leftmargin=2em, itemsep=0.2em]
\item The cut \(P\) is \emph{source-faithful for \(\Omega\)} if
      \[
      \mathfrak S_{\Omega}(P)=0.
      \]
\item The pair \((A,P)\) is \emph{one-step faithful for \(\Omega\)} if
      \[
      \mathfrak D_{\Omega}(P,A)=0.
      \]
\item The pair \((A,P)\) is \emph{task-exact for \(\Omega\)} if it is
      source-faithful, one-step faithful, and task-perfect.
\end{enumerate}
\end{definition}

\begin{remark}[What the defect package separates]
\label{rem:task-visible-defect-separation}\lean{CoherenceCalculus.sourceVisibleLoss, CoherenceCalculus.directObservableTransfer, CoherenceCalculus.fixedSourceObservableSchurReturnTransfer, CoherenceCalculus.taskVisibleDefectPackage}
These three channels measure different failures. Source-visible loss measures
what the cut removes before any dynamics act. Direct observable transfer
measures the one-step influence of the hidden source component. The fixed-source
Schur return transfer measures the hidden residual that re-enters after exact
elimination. The point of \(\mathcal E_{\Omega}(P,A;z)\) is that those failure
modes are no longer conflated.
\end{remark}

\subsection{Exact and approximate correction dimension}

\begin{theorem}[Exact correction factorization criterion]
\label{thm:exact-correction-factorization}\lean{CoherenceCalculus.ExactCorrectionFactorizationData, CoherenceCalculus.exact_correction_factorization}
Let \(T:S\to V\) and \(C:S\to R\) be linear maps. There exists a unique linear
map
\[
D_C:\Img(C)\to V
\]
with
\[
T=D_C\circ C
\]
if and only if
\[
\Ker(C)\subseteq \Ker(T).
\]
In finite dimensions, \(D_C\) extends to a linear decoder \(D:R\to V\) if one
wants a map on the whole correction space.
\end{theorem}

\begin{proof}
If \(T=D_C\circ C\) and \(Cx=0\), then \(Tx=D_C(Cx)=0\), so
\(\Ker(C)\subseteq\Ker(T)\).

Conversely, assume \(\Ker(C)\subseteq\Ker(T)\). Define \(D_C\) on \(\Img(C)\) by
\[
D_C(Cx):=Tx.
\]
This is well defined: if \(Cx=Cx'\), then \(C(x-x')=0\), so
\(x-x'\in\Ker(C)\subseteq\Ker(T)\), hence \(Tx=Tx'\).
If \(y\in\Img(C)\), write \(y=Cx\). Any linear map \(E:\Img(C)\to V\) with
\(T=E\circ C\) must satisfy
\[
E(y)=E(Cx)=Tx=D_C(Cx)=D_C(y),
\]
so \(E=D_C\). This proves uniqueness.
\end{proof}

\begin{remark}[Quotient-free exact correction]
\label{rem:quotient-free-exact-correction}\lean{CoherenceCalculus.ExactCorrectionFactorizationData, CoherenceCalculus.exact_correction_factorization}
The theorem is stated on \(\Img(C)\) rather than on the abstract quotient
\(S/\Ker(C)\). This is the exact same mathematical content, but it keeps the
factorization in an explicit image space and avoids introducing a quotient as a
primitive object.
\end{remark}

\begin{corollary}[Minimal exact correction dimension]
\label{cor:minimal-exact-correction-dimension}\lean{CoherenceCalculus.ObservableClosureRankData, CoherenceCalculus.observableClosureRank, CoherenceCalculus.minimal_exact_correction_dimension}
Let \(T:S\to V\) be linear with \(S\) and \(V\) finite-dimensional. The least
integer \(m\) for which there exist a space \(R\) with \(\dim R=m\) and maps
\[
C:S\to R,
\qquad
D:R\to V
\]
satisfying \(T=D\circ C\) is \(\rank(T)\).
\end{corollary}

\begin{proof}
Take \(R=\Img(T)\), let \(C:=T\), and let \(D:\Img(T)\hookrightarrow V\) be the
inclusion. Then \(T=D\circ C\) and \(\dim R=\rank(T)\).

Conversely, if \(T=D\circ C\), then
\[
\rank(T)\le \rank(C)\le \dim R.
\]
So every exact correction space has dimension at least \(\rank(T)\).
\end{proof}

\begin{definition}[Observable closure rank]
\label{def:observable-closure-rank}\lean{CoherenceCalculus.ObservableClosureRankData, CoherenceCalculus.observableClosureRank}
For a linear transfer \(T\), define its \emph{observable closure rank} by
\[
\rho(T):=\rank(T).
\]
If \(\Omega=(S,V,O)\) is an observable package in the Schur setting, define
\[
\rho_{\Omega}(P,A;z):=
\rank\bigl(\widetilde{\mathfrak R}_{\Omega}(P,A;z)\bigr).
\]
By Corollary~\ref{cor:minimal-exact-correction-dimension}, this is exactly the
minimal number of linear returned coordinates needed for exact closure of the
selected workload. If \(P(S)\subseteq S\), then
Proposition~\ref{prop:fixed-source-return-restricts} shows that this is also
the rank of the visible-workload return transfer
\(\mathfrak R_{\Omega}(P,A;z)\).
\end{definition}

\begin{corollary}[Zero rank is task-perfect coherence]
\label{cor:zero-closure-rank}\lean{CoherenceCalculus.zero_observable_closure_rank}
For every observable package \(\Omega=(S,V,O)\),
\[
\rho_{\Omega}(P,A;z)=0
\quad\Longleftrightarrow\quad
\widetilde{\mathfrak R}_{\Omega}(P,A;z)=0.
\]
\end{corollary}

\begin{proof}
A linear map has rank zero if and only if it is the zero map.
\end{proof}

\begin{definition}[Best \(m\)-dimensional correction error]
\label{def:best-m-correction-error}\lean{CoherenceCalculus.BestCorrectionErrorProfile, CoherenceCalculus.bestCorrectionError}
Let \(S\) and \(V\) be finite-dimensional normed spaces, let \(T:S\to V\) be
linear, and let \(m\in\bbN\). Define
\[
e_m(T):=
\inf_{\dim R\le m}
\inf_{C:S\to R,\ D:R\to V}
\|T-D\circ C\|_{\mathrm{op}}.
\]
\end{definition}

\begin{proposition}[Rank-constrained form and continuity of \(e_m\)]
\label{prop:best-m-correction-distance}\lean{CoherenceCalculus.CorrectionDistanceData, CoherenceCalculus.best_m_correction_distance}
Let \(S\) and \(V\) be finite-dimensional normed spaces. For every linear map
\(T:S\to V\),
\[
e_m(T)
=
\inf\{\|T-F\|_{\mathrm{op}}:\rank(F)\le m\}.
\]
Moreover, for all linear maps \(T,T':S\to V\),
\[
|e_m(T)-e_m(T')|
\le
\|T-T'\|_{\mathrm{op}}.
\]
\end{proposition}

\begin{proof}
If \(F=D\circ C\) with \(\dim R\le m\), then
\[
\rank(F)\le \rank(C)\le \dim R\le m.
\]
So every admissible factorization contributes a rank-at-most-\(m\) operator,
which proves
\[
e_m(T)\ge \inf\{\|T-F\|_{\mathrm{op}}:\rank(F)\le m\}.
\]

Conversely, if \(F:S\to V\) has \(\rank(F)\le m\), then \(F\) factors through
the \(m\)-dimensional space \(\Img(F)\): let
\[
C_F:S\to \Img(F),\qquad C_Fx:=Fx,
\]
and let \(D_F:\Img(F)\hookrightarrow V\) be the inclusion. Then
\[
F=D_F\circ C_F,
\]
so \(F\) is admissible in Definition~\ref{def:best-m-correction-error}. This
proves the rank-constrained formula.

For the Lipschitz bound, fix any rank-at-most-\(m\) operator \(F\). Then
\[
\|T-F\|_{\mathrm{op}}
\le
\|T-T'\|_{\mathrm{op}}+\|T'-F\|_{\mathrm{op}}.
\]
Taking the infimum over such \(F\) gives
\[
e_m(T)\le \|T-T'\|_{\mathrm{op}}+e_m(T').
\]
Interchanging \(T\) and \(T'\) yields the reverse inequality and hence the
claim.
\end{proof}

\subsection{Structured correction classes}

\begin{definition}[Structured correction class]
\label{def:structured-correction-class}\lean{CoherenceCalculus.StructuredCorrectionClass}
Let \(S\) and \(V\) be finite-dimensional normed spaces. For \(m\ge 0\), a
\emph{structured \(m\)-channel correction datum} from \(S\) to \(V\) is a
triple
\[
(R,C,D)
\]
where \(R\) is a vector space with \(\dim R\le m\), \(C:S\to R\) is linear,
and \(D:R\to V\) is linear. A \emph{structured correction class}
\(\mathcal C\) assigns to each \(m\ge 0\) a chosen subset
\[
\mathcal C_m(S,V)
\]
of such data.
\end{definition}

\begin{definition}[Structured closure width and structured closure rank]
\label{def:structured-closure-width}\lean{CoherenceCalculus.StructuredClosureProfile, CoherenceCalculus.structuredClosureWidth, CoherenceCalculus.structuredClosureRank}
Let \(\mathcal C\) be a structured correction class and let \(T:S\to V\) be
linear. Define the \emph{structured \(m\)-channel closure width} by
\[
e_m^{\mathcal C}(T)
:=
\inf_{(R,C,D)\in \mathcal C_m(S,V)}
\|T-D\circ C\|_{\mathrm{op}}.
\]
If there exists some \(m\) and some \((R,C,D)\in \mathcal C_m(S,V)\) with
\(T=D\circ C\), define the \emph{structured closure rank} \(\rho^{\mathcal C}(T)\)
to be the least such \(m\).
\end{definition}

\begin{proposition}[Universal lower envelope]
\label{prop:structured-width-lower-envelope}\lean{CoherenceCalculus.structured_width_lower_envelope}
Let \(\mathcal C\) be a structured correction class and let \(T:S\to V\) be
linear. Then
\[
e_m(T)\le e_m^{\mathcal C}(T)
\qquad (m\ge 0).
\]
If \(\rho^{\mathcal C}(T)\) is defined, then
\[
\rho(T)\le \rho^{\mathcal C}(T).
\]
\end{proposition}

\begin{proof}
Every datum \((R,C,D)\in \mathcal C_m(S,V)\) is in particular an admissible
\(m\)-channel correction datum in the sense of
Definition~\ref{def:best-m-correction-error}. Therefore the infimum over all
admissible data is bounded above by the infimum over the smaller structured
class.

If \(\rho^{\mathcal C}(T)=m\), then there exists
\((R,C,D)\in \mathcal C_m(S,V)\) with \(T=D\circ C\). Hence \(T\) admits an
exact \(m\)-channel factorization, so Corollary~\ref{cor:minimal-exact-correction-dimension}
gives \(\rho(T)\le m\).
\end{proof}

\begin{proposition}[Monotonicity for nested structured classes]
\label{prop:structured-width-monotonicity}\lean{CoherenceCalculus.structured_width_monotonicity}
Assume \(\mathcal C_m(S,V)\subseteq \mathcal C_{m+1}(S,V)\) for every \(m\).
Then
\[
e_{m+1}^{\mathcal C}(T)\le e_m^{\mathcal C}(T)
\qquad (m\ge 0).
\]
\end{proposition}

\begin{proof}
The infimum over the larger admissible set \(\mathcal C_{m+1}(S,V)\) cannot
exceed the infimum over the smaller set \(\mathcal C_m(S,V)\).
\end{proof}

\begin{remark}[Native examples]
\label{rem:structured-correction-native-examples}\lean{CoherenceCalculus.StructuredCorrectionClass, CoherenceCalculus.StructuredClosureProfile}
The class \(\mathcal C\) is where genuinely realizable restrictions enter:
equivariance, tower transport order, locality, block sparsity, or other
manuscript-native admissibility conditions. The unstructured width \(e_m\) is
the universal lower envelope; \(e_m^{\mathcal C}\) measures what remains once a
realizable design family is fixed.
\end{remark}

\subsection{Canonical optimal approximate correction}

\begin{definition}[Singular values]
\label{def:singular-values}\lean{CoherenceCalculus.SingularValueProfile, CoherenceCalculus.singularValue}
Let \(S\) and \(V\) be finite-dimensional inner-product spaces, and let
\(T:S\to V\) be linear. The \emph{singular values} of \(T\) are the nonnegative
numbers
\[
\sigma_1(T)\ge \sigma_2(T)\ge \cdots \ge \sigma_{\dim S}(T)\ge 0
\]
obtained by taking square roots of the eigenvalues of the positive
semidefinite operator \(T^\ast T\), repeated with multiplicity and listed in
nonincreasing order.
\end{definition}

\begin{theorem}[Optimal \(m\)-channel approximation]
\label{thm:optimal-m-channel-approximation}\lean{CoherenceCalculus.optimal_m_channel_approximation}
Let \(S\) and \(V\) be finite-dimensional inner-product spaces, let
\(T:S\to V\) be linear, and let
\[
\sigma_1(T)\ge \cdots \ge \sigma_{\dim S}(T)\ge 0
\]
be its singular values. Then there exist orthonormal vectors
\(u_1,\dots,u_{\dim S}\in S\) and an orthonormal family
\(v_1,\dots,v_r\in V\), where \(r=\rank(T)\), such that
\[
Tx=\sum_{i=1}^{r}\sigma_i(T)\,\langle x,u_i\rangle\,v_i
\qquad (x\in S).
\]
For \(m\ge 0\), define
\[
T^{[m]}x
:=
\sum_{i=1}^{\min(m,r)} \sigma_i(T)\,\langle x,u_i\rangle\,v_i.
\]
Then:
\begin{enumerate}[label=(\roman*), leftmargin=2em, itemsep=0.2em]
\item \(\rank(T^{[m]})\le m\);
\item \(e_m(T)=\|T-T^{[m]}\|_{\mathrm{op}}=\sigma_{m+1}(T)\), with the
      convention \(\sigma_{m+1}(T)=0\) when \(m\ge r\);
\item if \(U_m:=\Span\{u_1,\dots,u_{\min(m,r)}\}\), the orthogonal projector
      \(C_{T,m}:S\to U_m\) and the map \(D_{T,m}:U_m\to V\) defined by
      \[
      D_{T,m}(u_i):=\sigma_i(T)v_i
      \qquad
      (1\le i\le \min(m,r))
      \]
      satisfy
      \[
      T^{[m]}=D_{T,m}\circ C_{T,m}.
      \]
\end{enumerate}
\end{theorem}

(Proof in Appendix~\ref{sec:appendix-elimination}.)

\begin{corollary}[Optimal approximate returned correction for a workload]
\label{cor:optimal-returned-correction-for-workload}\lean{CoherenceCalculus.ObservableOptimalCorrectionData, CoherenceCalculus.optimal_returned_correction_for_workload}
Let \(\Omega=(S,V,O)\) be a finite-dimensional observable package on \(\cH\),
let \(P\) be an orthogonal projection, and assume
\(\widetilde{\mathfrak R}_{\Omega}(P,A;z)\) is defined. Then
\[
e_m\!\bigl(\widetilde{\mathfrak R}_{\Omega}(P,A;z)\bigr)
=
\sigma_{m+1}\!\bigl(\widetilde{\mathfrak R}_{\Omega}(P,A;z)\bigr).
\]
Accordingly, the visible workload admits an explicit optimal \(m\)-channel
returned correction obtained by truncating the singular expansion of the
observable Schur return transfer.
\end{corollary}

\begin{proof}
Apply Theorem~\ref{thm:optimal-m-channel-approximation} to the linear map
\[
T:=\widetilde{\mathfrak R}_{\Omega}(P,A;z):S\to V.
\]
\end{proof}

\begin{remark}[Canonicity and spectral gaps]
\label{rem:optimal-correction-gap}\lean{CoherenceCalculus.SingularValueProfile, CoherenceCalculus.optimal_m_channel_approximation}
If \(\sigma_m(T)>\sigma_{m+1}(T)\), then the leading singular subspace
\[
U_m=\Span\{u_i:\sigma_i(T)>\sigma_{m+1}(T)\}
\]
is uniquely determined by \(T^\ast T\), so the source projector \(C_{T,m}\) in
Theorem~\ref{thm:optimal-m-channel-approximation} is canonical. If
\(\sigma_m(T)=\sigma_{m+1}(T)\), then optimal \(m\)-channel corrections need
not be unique; the correct output is an optimal class rather than a singleton.
\end{remark}

\begin{remark}[Historical origin]
\label{rem:eckart-young-origin}\lean{CoherenceCalculus.optimal_m_channel_approximation}
The optimal low-rank approximation theorem is historically rooted in the work
of Eckart and Young, with the broader unitarily invariant norm perspective due
to Mirsky \cite{EckartYoung1936,Mirsky1960}. The proof used here is included
explicitly because later formal coverage must match the actual argument rather
than a citation boundary.
\end{remark}

\subsection{Design transforms and covariance}

\begin{definition}[Unitary design transform of a workload datum]
\label{def:observable-design-transform}\lean{CoherenceCalculus.ObservableDesignTransform}
Let \(U:\cH\to\cH\) be unitary. Define the transformed
operator, cut, and observable package by
\[
A^U:=UAU^{-1},
\qquad
P^U:=UPU^{-1},
\qquad
\Omega^U:=(US,V,OU^{-1}).
\]
\end{definition}

\begin{theorem}[Task-visible covariance under unitary design transforms]
\label{thm:observable-transfer-covariance}\lean{CoherenceCalculus.ObservableTransferCovarianceData, CoherenceCalculus.observable_transfer_covariance}
Let \(U:\cH\to\cH\) be unitary and write
\(\Omega=(S,V,O)\), \(\Omega^U=(US,V,OU^{-1})\),
\(A^U:=UAU^{-1}\), and \(P^U:=UPU^{-1}\). If
\(\CCSigmaP{P}{z}[A]\) is defined, then \(\CCSigmaP{P^U}{z}[A^U]\) is defined,
and
\[
\CCSigmaP{P^U}{z}[A^U]
=
U\,\CCSigmaP{P}{z}[A]\,U^{-1}.
\]
Moreover,
\[
\mathfrak S_{\Omega^U}(P^U)
=
\mathfrak S_{\Omega}(P)\circ U^{-1}\big|_{US},
\]
\[
\mathfrak D_{\Omega^U}(P^U,A^U)
=
\mathfrak D_{\Omega}(P,A)\circ U^{-1}\big|_{US},
\]
and
\[
\widetilde{\mathfrak R}_{\Omega^U}(P^U,A^U;z)
=
\widetilde{\mathfrak R}_{\Omega}(P,A;z)\circ U^{-1}\big|_{US}.
\]
\end{theorem}

\begin{proof}
Let \(Q:=I-P\). Then
\[
I-P^U
=
I-UPU^{-1}
=
U(I-P)U^{-1}
=
UQU^{-1}.
\]
Hence
\[
(I-P^U)A^U(I-P^U)
=
UQAQU^{-1}.
\]
Therefore
\[
I-z(I-P^U)A^U(I-P^U)
=
U(I-zQAQ)U^{-1},
\]
so invertibility of \(I-zQAQ\) implies invertibility of
\(I-z(I-P^U)A^U(I-P^U)\), with inverse
\[
\bigl(I-z(I-P^U)A^U(I-P^U)\bigr)^{-1}
=
U(I-zQAQ)^{-1}U^{-1}.
\]
Substituting this into the Schur formula gives
\begin{align*}
\CCSigmaP{P^U}{z}[A^U]
&=
P^U A^U (I-P^U)
\bigl(I-z(I-P^U)A^U(I-P^U)\bigr)^{-1}
(I-P^U)A^U P^U \\
&=
UPAQ(I-zQAQ)^{-1}QAPU^{-1} \\
&=
U\,\CCSigmaP{P}{z}[A]\,U^{-1}.
\end{align*}

For the source-visible loss,
\begin{align*}
\mathfrak S_{\Omega^U}(P^U)
&=
OU^{-1}(I-P^U)\big|_{US} \\
&=
OU^{-1}U(I-P)U^{-1}\big|_{US} \\
&=
O(I-P)U^{-1}\big|_{US}
=
\mathfrak S_{\Omega}(P)\circ U^{-1}\big|_{US}.
\end{align*}
The direct-transfer identity is the same computation with \(AQ\) in place of
\(I-P\):
\begin{align*}
\mathfrak D_{\Omega^U}(P^U,A^U)
&=
OU^{-1}A^U(I-P^U)\big|_{US} \\
&=
OAQ\,U^{-1}\big|_{US}
=
\mathfrak D_{\Omega}(P,A)\circ U^{-1}\big|_{US}.
\end{align*}
Finally,
\begin{align*}
\widetilde{\mathfrak R}_{\Omega^U}(P^U,A^U;z)
&=
OU^{-1}\,\CCSigmaP{P^U}{z}[A^U]\,P^U\big|_{US} \\
&=
OU^{-1}U\,\CCSigmaP{P}{z}[A]\,U^{-1}UPU^{-1}\big|_{US} \\
&=
O\,\CCSigmaP{P}{z}[A]\,P\,U^{-1}\big|_{US} \\
&=
\widetilde{\mathfrak R}_{\Omega}(P,A;z)\circ U^{-1}\big|_{US}.
\end{align*}
\end{proof}

\begin{proposition}[Source-isomorphism bounds for closure widths]
\label{prop:source-isomorphism-width-bounds}\lean{CoherenceCalculus.SourceIsomorphismWidthData, CoherenceCalculus.source_isomorphism_width_bounds}
Let \(W:S'\to S\) be an isomorphism between finite-dimensional normed spaces,
and let \(T:S\to V\) be linear. Then
\[
\|W^{-1}\|_{\mathrm{op}}^{-1}e_m(T)
\le
e_m(T\circ W)
\le
\|W\|_{\mathrm{op}}e_m(T).
\]
\end{proposition}

\begin{proof}
If \(F:S\to V\) has \(\rank(F)\le m\), then \(F\circ W:S'\to V\) also has rank
at most \(m\). Therefore
\[
e_m(T\circ W)
\le
\|T\circ W-F\circ W\|_{\mathrm{op}}
\le
\|T-F\|_{\mathrm{op}}\|W\|_{\mathrm{op}}.
\]
Taking the infimum over rank-at-most-\(m\) maps \(F\) gives the upper bound.

Now apply the same estimate to the pair \((T\circ W,W^{-1})\):
\[
e_m(T)
=
e_m\bigl((T\circ W)\circ W^{-1}\bigr)
\le
\|W^{-1}\|_{\mathrm{op}}\,e_m(T\circ W).
\]
Rearranging yields the lower bound.
\end{proof}

\begin{corollary}[Unitary invariance of workload closure data]
\label{cor:unitary-workload-invariance}\lean{CoherenceCalculus.UnitaryWorkloadInvarianceData, CoherenceCalculus.unitary_workload_invariance}
In the setting of Theorem~\ref{thm:observable-transfer-covariance}, assume
\(U\) is unitary. Then
\[
\rho_{\Omega^U}(P^U,A^U;z)=\rho_{\Omega}(P,A;z),
\]
\[
e_m\!\bigl(\widetilde{\mathfrak R}_{\Omega^U}(P^U,A^U;z)\bigr)
=
e_m\!\bigl(\widetilde{\mathfrak R}_{\Omega}(P,A;z)\bigr),
\]
and
\[
\sigma_i\!\bigl(\widetilde{\mathfrak R}_{\Omega^U}(P^U,A^U;z)\bigr)
=
\sigma_i\!\bigl(\widetilde{\mathfrak R}_{\Omega}(P,A;z)\bigr)
\qquad (1\le i\le \dim S).
\]
\end{corollary}

\begin{proof}
Let \(W:=U^{-1}\big|_{US}:US\to S\). Because \(U\) is unitary, \(W\) is
unitary as well, so \(\|W\|_{\mathrm{op}}=\|W^{-1}\|_{\mathrm{op}}=1\). The
covariance theorem gives
\[
\widetilde{\mathfrak R}_{\Omega^U}(P^U,A^U;z)
=
\widetilde{\mathfrak R}_{\Omega}(P,A;z)\circ W.
\]
Composition with the isomorphism \(W\) does not change rank, so the closure
rank is unchanged. Proposition~\ref{prop:source-isomorphism-width-bounds}
therefore gives equality of the widths \(e_m\). Finally,
Theorem~\ref{thm:optimal-m-channel-approximation} identifies
\(\sigma_i(T)\) with \(e_{i-1}(T)\), so equality of all widths implies
equality of all singular values.
\end{proof}

\subsection{Observable closure Gramians and invariants}

\begin{definition}[Observable closure Gramian]
\label{def:observable-closure-gramian}\lean{CoherenceCalculus.ObservableClosureGramianData, CoherenceCalculus.observableClosureGramian}
Let \(\Omega=(S,V,O)\) be a finite-dimensional observable package on \(\cH\),
and assume \(\widetilde{\mathfrak R}_{\Omega}(P,A;z)\) is defined. Set
\[
T_{\Omega,P,z}
:=
\widetilde{\mathfrak R}_{\Omega}(P,A;z):S\to V.
\]
The \emph{observable closure Gramian} is
\[
G_{\Omega}(P,A;z):=T_{\Omega,P,z}^\ast T_{\Omega,P,z}:S\to S.
\]
\end{definition}

\begin{definition}[Observable closure invariant family]
\label{def:observable-closure-invariants}\lean{CoherenceCalculus.ObservableClosureInvariantFamily, CoherenceCalculus.observableClosureInvariantFamilyOf}
In the setting of Definition~\ref{def:observable-closure-gramian}, define:
\begin{enumerate}[label=(\roman*), itemsep=0.2em]
\item the multiset
      \[
      \Lambda_{\Omega}(P,A;z):=
      \mathrm{spec}\bigl(G_{\Omega}(P,A;z)\bigr)
      \]
      of nonnegative eigenvalues of the closure Gramian;
\item the spectral moments
      \[
      M_k^{\Omega}(P,A;z):=
      \Tr\bigl(G_{\Omega}(P,A;z)^k\bigr)
      \qquad (k\ge 1);
      \]
\item the determinant polynomial
      \[
      \Delta_{\Omega}(P,A;z;t):=
      \det\bigl(I+tG_{\Omega}(P,A;z)\bigr).
      \]
\end{enumerate}
\end{definition}

\begin{theorem}[Finite observable closure package]
\label{thm:observable-closure-invariant-package}\lean{CoherenceCalculus.observable_closure_invariant_package}
Assume the setting of Definition~\ref{def:observable-closure-invariants}, and
let \(n:=\dim S\). Then:
\begin{enumerate}[label=(\roman*), itemsep=0.2em]
\item every element of \(\Lambda_{\Omega}(P,A;z)\) is nonnegative;
\item the following are equivalent:
      \[
      \Lambda_{\Omega}(P,A;z)=\{0,\dots,0\},
      \qquad
      G_{\Omega}(P,A;z)=0,
      \qquad
      \widetilde{\mathfrak R}_{\Omega}(P,A;z)=0;
      \]
\item if
      \[
      \Lambda_{\Omega}(P,A;z)=\{\lambda_1,\dots,\lambda_n\}
      \]
      with multiplicity in nonincreasing order, then
      \[
      \sigma_i\!\bigl(\widetilde{\mathfrak R}_{\Omega}(P,A;z)\bigr)
      =
      \sqrt{\lambda_i}
      \qquad (1\le i\le n),
      \]
      \[
      \rho_{\Omega}(P,A;z)=\#\{i:\lambda_i>0\},
      \qquad
      e_m\!\bigl(\widetilde{\mathfrak R}_{\Omega}(P,A;z)\bigr)
      =
      \sqrt{\lambda_{m+1}},
      \]
      with the convention \(\lambda_{m+1}=0\) for \(m\ge \rho_{\Omega}(P,A;z)\);
\item
      \[
      M_k^{\Omega}(P,A;z)=\sum_{i=1}^{n}\lambda_i^k,
      \qquad
      \Delta_{\Omega}(P,A;z;t)=\prod_{i=1}^{n}(1+t\lambda_i);
      \]
\item if \(U\) is unitary, then the transformed datum of
      Definition~\ref{def:observable-design-transform} has the same invariant
      family:
      \[
      \Lambda_{\Omega^U}(P^U,A^U;z)=\Lambda_{\Omega}(P,A;z),
      \]
      \[
      M_k^{\Omega^U}(P^U,A^U;z)=M_k^{\Omega}(P,A;z),
      \qquad
      \Delta_{\Omega^U}(P^U,A^U;z;t)=\Delta_{\Omega}(P,A;z;t).
      \]
\end{enumerate}
\end{theorem}

(Proof in Appendix~\ref{sec:appendix-elimination}.)

\begin{proposition}[Singular-value stability]
\label{prop:singular-value-stability}\lean{CoherenceCalculus.SingularValueStabilityData, CoherenceCalculus.singular_value_stability}
Let \(S\) and \(V\) be finite-dimensional inner-product spaces, and let
\(T,T':S\to V\) be linear. Then
\[
\bigl|\sigma_i(T)-\sigma_i(T')\bigr|
\le
\|T-T'\|_{\mathrm{op}}
\qquad (1\le i\le \dim S).
\]
\end{proposition}

\begin{proof}
By Theorem~\ref{thm:optimal-m-channel-approximation},
\[
\sigma_i(T)=e_{i-1}(T),
\qquad
\sigma_i(T')=e_{i-1}(T').
\]
Applying Proposition~\ref{prop:best-m-correction-distance} with \(m=i-1\)
gives the claim.
\end{proof}

\begin{corollary}[Continuity of the workload closure spectrum]
\label{cor:workload-closure-spectrum-continuity}\lean{CoherenceCalculus.WorkloadClosureSpectrumContinuityData, CoherenceCalculus.workload_closure_spectrum_continuity}
In the setting of Proposition~\ref{prop:fixed-source-return-continuity}, every
singular value
\[
P\longmapsto \sigma_i\!\bigl(\widetilde{\mathfrak R}_{\Omega}(P,A;z)\bigr)
\]
is continuous on the admissible cut family.
\end{corollary}

\begin{proof}
Combine Proposition~\ref{prop:fixed-source-return-continuity} with
Proposition~\ref{prop:singular-value-stability}.
\end{proof}

\subsection{Effective and null returned directions}

\begin{definition}[Observable null and effective projectors]
\label{def:observable-null-effective-projectors}\lean{CoherenceCalculus.ObservableEffectiveNullProjectors}
Let \(T:S\to V\) be linear between finite-dimensional inner-product spaces. The
\emph{observable null projector} and \emph{observable effective projector} are
\[
P_T^{\mathrm{null}}:=I-T^+T,
\qquad
P_T^{\mathrm{eff}}:=T^+T,
\]
where \(T^+\) is the Moore--Penrose pseudoinverse.
\end{definition}

\begin{proposition}[Effective/null decomposition of a returned residual]
\label{prop:observable-effective-null-decomposition}\lean{CoherenceCalculus.observable_effective_null_decomposition}
In the setting of
Definition~\ref{def:observable-null-effective-projectors},
\(P_T^{\mathrm{null}}\) is the orthogonal projector onto \(\Ker(T)\),
\(P_T^{\mathrm{eff}}\) is the orthogonal projector onto \((\Ker T)^\perp\), and
every \(r\in S\) decomposes as
\[
r=r_{\mathrm{eff}}+r_{\mathrm{null}},
\qquad
r_{\mathrm{eff}}:=P_T^{\mathrm{eff}}r,
\qquad
r_{\mathrm{null}}:=P_T^{\mathrm{null}}r,
\]
with
\[
T(r_{\mathrm{null}})=0,
\qquad
T(r)=T(r_{\mathrm{eff}}).
\]
\end{proposition}

\begin{proof}
Apply Lemma~\ref{lem:trace-ker} to the constraint map \(C:=T\) in
Appendix~\ref{sec:boundary-defect}. It shows that
\[
I-T^+T=P_T^{\mathrm{null}}
\]
is the orthogonal projector onto \(\Ker(T)\). Therefore
\[
P_T^{\mathrm{eff}}=I-P_T^{\mathrm{null}}
\]
is the orthogonal projector onto \((\Ker T)^\perp\). The displayed
decomposition is just
\[
r=P_T^{\mathrm{eff}}r+P_T^{\mathrm{null}}r.
\]
Since \(r_{\mathrm{null}}\in\Ker(T)\), we have \(T(r_{\mathrm{null}})=0\), and
hence
\[
T(r)=T(r_{\mathrm{eff}})+T(r_{\mathrm{null}})=T(r_{\mathrm{eff}}).
\]
\end{proof}

\begin{quote}\small
\textbf{Interpretation.}
For a selected workload, \(r_{\mathrm{null}}\) is the shadow component of a
returned residual that the workload never sees, while \(r_{\mathrm{eff}}\) is
the minimal Hilbert representative that carries exactly the observable effect.
\end{quote}

\subsection{Observable Schur jets}

\begin{definition}[Observable Schur-jet truncation]
\label{def:observable-schur-jet-truncation}\lean{CoherenceCalculus.ObservableSchurJetData, CoherenceCalculus.observableSchurJetTruncation}
Let \(\Omega=(S,V,O)\) be an observable package on \(\cH\), let \(P\) be an
orthogonal projection with \(Q:=I-P\), and assume \(\|zQAQ\|<1\). For
\(N\ge 1\), define the truncated observable return transfer
\[
\widetilde{\mathfrak R}_{\Omega}^{[N]}(P,A;z)
:=
O\left(
\sum_{n=0}^{N-1} z^n\,PAQ(QAQ)^nQAP
\right)P\Big|_{S}.
\]
\end{definition}

\begin{theorem}[Observable Schur-jet remainder]
\label{thm:observable-schur-jet-remainder}\lean{CoherenceCalculus.observable_schur_jet_remainder}
In the setting of Definition~\ref{def:observable-schur-jet-truncation},
\[
\widetilde{\mathfrak R}_{\Omega}(P,A;z)
-\widetilde{\mathfrak R}_{\Omega}^{[N]}(P,A;z)
=
O\Bigl(
z^N PAQ(QAQ)^N(I-zQAQ)^{-1}QAP
\Bigr)P\Big|_{S}.
\]
Consequently,
\[
\bigl\|
\widetilde{\mathfrak R}_{\Omega}(P,A;z)
-\widetilde{\mathfrak R}_{\Omega}^{[N]}(P,A;z)
\bigr\|_{\mathrm{op}}
\le
\|O\|\cdot \|PAQ\|\cdot \|zQAQ\|^N\cdot \|(I-zQAQ)^{-1}\|\cdot \|QAP\|.
\]
\end{theorem}

(Proof in Appendix~\ref{sec:appendix-elimination}.)

\begin{definition}[Observable closure order]
\label{def:observable-closure-order}\lean{CoherenceCalculus.ObservableClosureOrderData, CoherenceCalculus.observableClosureOrder}
In the setting of Definition~\ref{def:observable-schur-jet-truncation}, the
\emph{exact observable closure order} is the least \(N\ge 1\) such that
\[
\widetilde{\mathfrak R}_{\Omega}(P,A;z)
=
\widetilde{\mathfrak R}_{\Omega}^{[N]}(P,A;z).
\]
For a tolerance \(\varepsilon>0\), the \emph{\(\varepsilon\)-observable closure
order} is the least \(N\ge 1\) with
\[
\bigl\|
\widetilde{\mathfrak R}_{\Omega}(P,A;z)
-\widetilde{\mathfrak R}_{\Omega}^{[N]}(P,A;z)
\bigr\|_{\mathrm{op}}
\le \varepsilon.
\]
\end{definition}

\subsection{Stochastic and coded closure}

\begin{quote}\small
\textbf{Boundary.}
The deterministic transfer \(T\) remains primary. The stochastic layer does not
replace it; it studies how \(T\) can be realized by kernels and finite-bit
codes, with explicit mean, covariance, rate, entropy, and seed accounting.
\end{quote}

\begin{definition}[Support and probabilistic workload data]
\label{def:stochastic-workload-data}\lean{CoherenceCalculus.SupportWorkloadData, CoherenceCalculus.ProbabilisticWorkloadData}
Let \(S\) be a finite-dimensional Hilbert space and let \(T:S\to V\) be
linear.
\begin{enumerate}[label=(\roman*), leftmargin=2em, itemsep=0.2em]
\item A \emph{support workload datum} for \(T\) is a Borel set
      \[
      K\subseteq S.
      \]
\item A \emph{probabilistic workload datum} for \(T\) is a Borel probability
      measure \(\mu\in\mathcal P(S)\) with finite second moment:
      \[
      \int_S \|x\|^2\,d\mu(x)<\infty.
      \]
\end{enumerate}
\end{definition}

\begin{definition}[Stochastic closure kernel]
\label{def:stochastic-closure-kernel}\lean{CoherenceCalculus.StochasticClosureKernel}
Let \(K\subseteq S\) be a Borel workload set. A \emph{stochastic closure
kernel} from \(K\) to \(V\) is a Markov kernel
\[
\kappa:K\rightsquigarrow V
\]
such that
\[
\int_V \|v\|^2\,\kappa_x(dv)<\infty
\qquad (\forall x\in K).
\]
Its \emph{mean map} is
\[
M_{\kappa}(x):=\int_V v\,\kappa_x(dv),
\]
and its \emph{covariance operator field} is
\[
\Gamma_{\kappa}(x)
:=
\int_V (v-M_{\kappa}(x))\otimes(v-M_{\kappa}(x))\,\kappa_x(dv),
\]
where
\[
(u\otimes w)z:=u\,\langle z,w\rangle.
\]
\end{definition}

\begin{definition}[Exactness and unbiasedness]
\label{def:stochastic-exactness-unbiasedness}\lean{CoherenceCalculus.DeterministicZeroNoiseData}
Let \(T:K\to V\) be measurable and let \(\kappa\) be a stochastic closure
kernel on \(K\).
\begin{enumerate}[label=(\roman*), leftmargin=2em, itemsep=0.2em]
\item \(\kappa\) is \emph{exact} for \(T\) if
      \[
      \kappa_x=\delta_{T(x)}
      \qquad (\forall x\in K).
      \]
\item \(\kappa\) is \emph{unbiased} for \(T\) if
      \[
      M_{\kappa}(x)=T(x)
      \qquad (\forall x\in K).
      \]
\end{enumerate}
\end{definition}

\begin{theorem}[Pointwise bias--variance identity]
\label{thm:pointwise-bias-variance}\lean{CoherenceCalculus.StochasticBiasVarianceData, CoherenceCalculus.pointwise_bias_variance}
Let \(T:K\to V\) be measurable and let \(\kappa\) be a stochastic closure
kernel on \(K\). Then for every \(x\in K\),
\[
\int_V \|v-T(x)\|^2\,\kappa_x(dv)
=
\|M_{\kappa}(x)-T(x)\|^2
+
\Tr\Gamma_{\kappa}(x).
\]
\end{theorem}

\begin{proof}
Fix \(x\in K\), and write \(m:=M_{\kappa}(x)\), \(t:=T(x)\). Then
\[
v-t=(v-m)+(m-t),
\]
so
\[
\|v-t\|^2
=
\|v-m\|^2
+
2\,\mathrm{Re}\langle v-m,\ m-t\rangle
+
\|m-t\|^2.
\]
Integrating against \(\kappa_x\), the cross term vanishes because
\[
\int_V (v-m)\,\kappa_x(dv)=0.
\]
Therefore
\[
\int_V \|v-t\|^2\,\kappa_x(dv)
=
\int_V \|v-m\|^2\,\kappa_x(dv)
+
\|m-t\|^2.
\]
Since \(V\) is finite-dimensional,
\[
\int_V \|v-m\|^2\,\kappa_x(dv)=\Tr\Gamma_{\kappa}(x),
\]
which proves the identity.
\end{proof}

\begin{corollary}[Integrated bias--variance identity]
\label{cor:integrated-bias-variance}\lean{CoherenceCalculus.IntegratedBiasVarianceData, CoherenceCalculus.integrated_bias_variance}
Let \(\mu\) be a probabilistic workload datum for \(T\), and let \(\kappa\) be
a stochastic closure kernel on \(\operatorname{supp}\mu\). Then
\[
\int_S\int_V \|v-T(x)\|^2\,\kappa_x(dv)\,d\mu(x)
=
\|M_{\kappa}-T\|_{L^2(\mu;V)}^2
+
\int_S \Tr\Gamma_{\kappa}(x)\,d\mu(x).
\]
\end{corollary}

\begin{proof}
Integrate Theorem~\ref{thm:pointwise-bias-variance} against \(\mu\).
\end{proof}

\begin{proposition}[Zero covariance criterion]
\label{prop:zero-covariance-criterion}\lean{CoherenceCalculus.ZeroCovarianceCriterionData, CoherenceCalculus.zero_covariance_criterion}
Let \(\kappa\) be a stochastic closure kernel on \(K\). For a fixed \(x\in K\),
the following are equivalent:
\begin{enumerate}[label=(\roman*), leftmargin=2em, itemsep=0.2em]
\item \(\Gamma_{\kappa}(x)=0\);
\item \(\kappa_x=\delta_{M_{\kappa}(x)}\).
\end{enumerate}
\end{proposition}

\begin{proof}
If \(\kappa_x=\delta_{M_{\kappa}(x)}\), then the covariance is visibly zero.

Conversely, if \(\Gamma_{\kappa}(x)=0\), then
\[
0=\Tr\Gamma_{\kappa}(x)
=
\int_V \|v-M_{\kappa}(x)\|^2\,\kappa_x(dv),
\]
so \(v=M_{\kappa}(x)\) holds \(\kappa_x\)-almost surely. Therefore
\(\kappa_x=\delta_{M_{\kappa}(x)}\).
\end{proof}

\begin{corollary}[Deterministic closure as the zero-noise edge]
\label{cor:deterministic-zero-noise-edge}\lean{CoherenceCalculus.DeterministicZeroNoiseData, CoherenceCalculus.deterministic_zero_noise_edge}
Let \(T:K\to V\) be measurable and let \(\kappa\) be a stochastic closure
kernel on \(K\). Then \(\kappa\) is exact for \(T\) if and only if it is
unbiased for \(T\) and \(\Gamma_{\kappa}(x)=0\) for every \(x\in K\).
\end{corollary}

\begin{proof}
If \(\kappa_x=\delta_{T(x)}\), then \(M_{\kappa}(x)=T(x)\) and
\(\Gamma_{\kappa}(x)=0\).

Conversely, assume unbiasedness and zero covariance. By
Proposition~\ref{prop:zero-covariance-criterion},
\[
\kappa_x=\delta_{M_{\kappa}(x)}=\delta_{T(x)}
\qquad (\forall x\in K).
\]
\end{proof}

\begin{definition}[Seeded coded stochastic closure datum]
\label{def:seeded-coded-stochastic-closure-datum}\lean{CoherenceCalculus.SeededClosureCode}
Let \(K\subseteq S\) be a Borel workload set. A \emph{seeded coded stochastic
closure datum} from \(K\) to \(V\) is a quintuple
\[
\mathbf C=(\Xi,\pi,\mathcal A,\ell,(E,D))
\]
consisting of:
\begin{enumerate}[label=(\roman*), leftmargin=2em, itemsep=0.2em]
\item a standard Borel probability space \((\Xi,\pi)\), called the
      \emph{shared seed space};
\item a countable alphabet \(\mathcal A\);
\item a length function \(\ell:\mathcal A\to \bbN\) satisfying Kraft's
      inequality
      \[
      \sum_{a\in\mathcal A} 2^{-\ell(a)}\le 1;
      \]
\item a measurable encoder kernel
      \[
      E:K\times\Xi\rightsquigarrow \mathcal A;
      \]
\item a measurable decoder
      \[
      D:\mathcal A\times\Xi\to V.
      \]
\end{enumerate}
For \(x\in K\), \(\xi\in\Xi\), and \(a\in\mathcal A\), define the decoded
random output
\[
\widehat T_{\mathbf C}(x,\xi,a):=D(a,\xi).
\]
The induced stochastic closure kernel is
\[
\kappa^{\mathbf C}_x(B)
:=
\int_{\Xi}\sum_{a\in\mathcal A}
E_{x,\xi}(a)\,\mathbf 1_B\!\bigl(D(a,\xi)\bigr)\,d\pi(\xi).
\]
Its mean map and covariance field are
\[
M_{\mathbf C}:=M_{\kappa^{\mathbf C}},
\qquad
\Gamma_{\mathbf C}:=\Gamma_{\kappa^{\mathbf C}}.
\]
Unless explicitly stated otherwise, the rate functionals below charge the
emitted symbols but not the shared seed.
\end{definition}

\begin{definition}[Rate profiles]
\label{def:stochastic-rate-profiles}\lean{CoherenceCalculus.ClosureRateProfile, CoherenceCalculus.pointwiseCodeLength, CoherenceCalculus.worstCaseCodeLength, CoherenceCalculus.meanCodeLength}
Let \(\mathbf C\) be a seeded coded stochastic closure datum from a support
workload \(K\subseteq S\) to \(V\). Its pointwise expected code length is
\[
L_{\mathbf C}(x)
:=
\int_{\Xi}\sum_{a\in\mathcal A} E_{x,\xi}(a)\,\ell(a)\,d\pi(\xi).
\]
For a support workload \(K\subseteq S\), define
\[
L_K(\mathbf C):=\sup_{x\in K} L_{\mathbf C}(x).
\]
For a probabilistic workload datum \(\mu\) with
\(\operatorname{supp}\mu\subseteq K\), define
\[
L_{\mu}(\mathbf C):=\int_S L_{\mathbf C}(x)\,d\mu(x).
\]
\end{definition}

\begin{definition}[Rate-constrained stochastic closure widths]
\label{def:stochastic-closure-widths}\lean{CoherenceCalculus.ClosureWidthProfile, CoherenceCalculus.worstCaseClosureWidth, CoherenceCalculus.meanSquareClosureWidth, CoherenceCalculus.worstCaseUnbiasedClosureWidth, CoherenceCalculus.meanSquareUnbiasedClosureWidth}
Let \(T:K\to V\) be measurable, and let \(\mathcal C\) be an admissible class
of seeded coded stochastic closure data on \(K\).
\begin{enumerate}[label=(\roman*), leftmargin=2em, itemsep=0.2em]
\item The \emph{worst-case RMS stochastic closure width} at budget \(b\) is
      \[
      \eta^{\mathcal C}_{b,K}(T)
      :=
      \inf_{\mathbf C\in\mathcal C,\ L_K(\mathbf C)\le b}
      \sup_{x\in K}
      \Bigl(
      \mathbb E\|\widehat T_{\mathbf C}(x,\xi,a)-T(x)\|^2
      \Bigr)^{1/2}.
      \]
\item For a probabilistic workload datum \(\mu\) with
      \(\operatorname{supp}\mu\subseteq K\), the
      \emph{mean-square stochastic closure width} is
      \[
      \eta^{\mathcal C}_{b,\mu}(T)
      :=
      \inf_{\mathbf C\in\mathcal C,\ L_{\mu}(\mathbf C)\le b}
      \left(
      \int_S
      \mathbb E\|\widehat T_{\mathbf C}(x,\xi,a)-T(x)\|^2\,d\mu(x)
      \right)^{1/2}.
      \]
\item The corresponding \emph{unbiased widths}
      \(u^{\mathcal C}_{b,K}(T)\) and \(u^{\mathcal C}_{b,\mu}(T)\) are defined
      by the same formulas with the additional condition \(M_{\mathbf C}=T\).
\end{enumerate}
\end{definition}

\begin{proposition}[Monotonicity of stochastic closure widths]
\label{prop:stochastic-width-monotonicity}\lean{CoherenceCalculus.StochasticWidthMonotonicityData, CoherenceCalculus.stochastic_width_monotonicity}
If \(b'\ge b\), then
\[
\eta^{\mathcal C}_{b',K}(T)\le \eta^{\mathcal C}_{b,K}(T),
\qquad
\eta^{\mathcal C}_{b',\mu}(T)\le \eta^{\mathcal C}_{b,\mu}(T),
\]
and likewise for the unbiased widths. If \(\mathcal C\subseteq \mathcal C'\),
then
\[
\eta^{\mathcal C'}_{b,*}(T)\le \eta^{\mathcal C}_{b,*}(T),
\qquad
u^{\mathcal C'}_{b,*}(T)\le u^{\mathcal C}_{b,*}(T).
\]
\end{proposition}

\begin{proof}
Each inequality follows because the relevant infimum is taken over a larger
admissible set.
\end{proof}

\begin{definition}[Unseeded codebook datum]
\label{def:unseeded-codebook-datum}\lean{CoherenceCalculus.UnseededCodebook}
An \emph{unseeded codebook datum} on a support workload \(K\subseteq S\)
consists of a finite alphabet
\[
\mathcal A=\{1,\dots,N\},
\]
decoder atoms \(v_1,\dots,v_N\in V\), and a measurable map
\[
\lambda:K\to \Delta_N,
\]
where
\[
\Delta_N:=\left\{(\lambda_1,\dots,\lambda_N)\in [0,1]^N:
\sum_{i=1}^N\lambda_i=1\right\}.
\]
Its induced stochastic closure kernel is
\[
\kappa^{\lambda}_x
=
\sum_{i=1}^{N}\lambda_i(x)\,\delta_{v_i}.
\]
\end{definition}

\begin{theorem}[Barycentric exactness criterion]
\label{thm:barycentric-exactness-criterion}\lean{CoherenceCalculus.BarycentricExactnessData, CoherenceCalculus.barycentric_exactness_criterion}
Let \(T:K\to V\) be measurable. The following are equivalent:
\begin{enumerate}[label=(\roman*), leftmargin=2em, itemsep=0.2em]
\item there exists an unseeded finite-alphabet unbiased code for \(T\) with
      decoder atoms \(v_1,\dots,v_N\);
\item there exists a measurable \(\lambda:K\to\Delta_N\) such that
      \[
      T(x)=\sum_{i=1}^{N}\lambda_i(x)\,v_i
      \qquad (\forall x\in K).
      \]
\end{enumerate}
\end{theorem}

\begin{proof}
If such a code exists, let \(\lambda_i(x)\) be the probability that the encoder
emits symbol \(i\) at source point \(x\). Then the decoded mean is
\[
M(x)=\sum_{i=1}^{N}\lambda_i(x)\,v_i.
\]
Unbiasedness gives \(M(x)=T(x)\).

Conversely, given such a barycentric selector \(\lambda\), let the encoder
sample \(i\in\{1,\dots,N\}\) according to \(\lambda(x)\), and let the decoder
output \(v_i\). Then
\[
\mathbb E[\widehat T(x)]
=
\sum_{i=1}^{N}\lambda_i(x)\,v_i
=
T(x),
\]
so the resulting code is unbiased.
\end{proof}

\begin{definition}[Closure hull rank]
\label{def:closure-hull-rank}\lean{CoherenceCalculus.ClosureHullRankData, CoherenceCalculus.closureHullRank}
For a measurable target \(T:K\to V\), define its \emph{closure hull rank}
\[
\operatorname{hrk}_{K}(T)
\]
to be the least \(N\in\bbN\cup\{\infty\}\) such that there exist decoder atoms
\[
v_1,\dots,v_N\in V
\]
and a measurable barycentric selector \(\lambda:K\to \Delta_N\) satisfying
\[
T(x)=\sum_{i=1}^{N}\lambda_i(x)\,v_i
\qquad (\forall x\in K).
\]
\end{definition}

\begin{corollary}[Exact fixed-length stochastic rate]
\label{cor:exact-fixed-length-stochastic-rate}\lean{CoherenceCalculus.FixedLengthRateData, CoherenceCalculus.exact_fixed_length_stochastic_rate}
If \(\operatorname{hrk}_{K}(T)=N<\infty\), then \(T\) admits an exact unbiased
fixed-length code using \(\lceil \log_2 N\rceil\) bits per source instance.
Conversely, no exact unbiased fixed-length code with fewer than
\(\lceil \log_2 N\rceil\) bits exists.
\end{corollary}

\begin{proof}
If \(\operatorname{hrk}_{K}(T)=N\), Theorem~\ref{thm:barycentric-exactness-criterion}
gives an exact unbiased code with \(N\) symbols, which can be transmitted with
\(\lceil \log_2 N\rceil\) bits.

Conversely, a fixed-length \(b\)-bit code has at most \(2^b\) symbols. If such
a code were exact and unbiased, Theorem~\ref{thm:barycentric-exactness-criterion}
would imply \(\operatorname{hrk}_{K}(T)\le 2^b\). Hence
\[
b\ge \lceil \log_2 \operatorname{hrk}_{K}(T)\rceil.
\]
\end{proof}

\begin{example}[Exact one-bit unbiased scalar closure]
\label{ex:one-bit-unbiased-scalar-closure}\lean{CoherenceCalculus.OneBitUnbiasedScalarClosureData, CoherenceCalculus.one_bit_unbiased_scalar_closure}
Take
\[
S=V=\bbR,
\qquad
K=[-1,1],
\qquad
T(x)=x.
\]
Let the alphabet be \(\mathcal A=\{-1,+1\}\), with both symbols assigned length
\(1\), and define
\[
\mathbb P(a=+1\mid x)=\frac{1+x}{2},
\qquad
\mathbb P(a=-1\mid x)=\frac{1-x}{2}.
\]
Decode by
\[
D(+1)=1,
\qquad
D(-1)=-1.
\]
Then
\[
\mathbb E[\widehat T(x)]
=
\frac{1+x}{2}(1)+\frac{1-x}{2}(-1)
=
x,
\]
so the code is exact in expectation using one bit. Its variance is
\[
\mathbb E\bigl[(\widehat T(x)-x)^2\bigr]=1-x^2.
\]
\end{example}

\begin{definition}[Closure entropy profile]
\label{def:closure-entropy-profile}\lean{CoherenceCalculus.ClosureEntropyProfile, CoherenceCalculus.closureEntropy}
Let \(\mathbf C\) be a seeded coded stochastic closure datum. Define
\[
h_{\mathbf C}(x)
:=
\int_{\Xi} H(E_{x,\xi})\,d\pi(\xi),
\]
where the Shannon entropy in bits is
\[
H(p):=-\sum_{a\in\mathcal A} p(a)\log_2 p(a),
\]
with the convention \(0\log_2 0:=0\).
\end{definition}

\begin{theorem}[Entropy lower bound]
\label{thm:entropy-lower-bound}\lean{CoherenceCalculus.ClosureEntropyProfile, CoherenceCalculus.entropy_lower_bound}
Let \(\mathbf C\) be a seeded coded stochastic closure datum from a support
workload \(K\subseteq S\) to \(V\). Then
\[
h_{\mathbf C}(x)\le L_{\mathbf C}(x)
\qquad (\forall x\in K).
\]
If \(\mu\) is a probabilistic workload datum with
\(\operatorname{supp}\mu\subseteq K\), then
\[
\int_S h_{\mathbf C}(x)\,d\mu(x)\le L_{\mu}(\mathbf C).
\]
\end{theorem}

\begin{proof}
Fix \(x\in K\) and \(\xi\in\Xi\), and write \(p_a:=E_{x,\xi}(a)\). Set
\[
c_{\ell}:=\sum_{a\in\mathcal A}2^{-\ell(a)}\le 1,
\qquad
q_a:=\frac{2^{-\ell(a)}}{c_{\ell}}.
\]
Then
\begin{align*}
\sum_{a\in\mathcal A} p_a\,\ell(a)
&=
-\sum_{a\in\mathcal A} p_a \log_2 q_a-\log_2 c_{\ell} \\
&=
H(p)+D_{\mathrm{KL}}(p\|q)-\log_2 c_{\ell} \\
&\ge H(p),
\end{align*}
because \(D_{\mathrm{KL}}(p\|q)\ge 0\) and \(-\log_2 c_{\ell}\ge 0\). Integrate
first over \(\xi\). This gives \(h_{\mathbf C}(x)\le L_{\mathbf C}(x)\) for
every \(x\in K\). If \(\mu\) is a probabilistic workload datum supported in
\(K\), integrating that pointwise inequality against \(\mu\) yields the second
claim.
\end{proof}

\begin{remark}[Historical origin]
\label{rem:stochastic-coding-origin}\lean{CoherenceCalculus.ClosureEntropyProfile, CoherenceCalculus.entropy_lower_bound}
The entropy lower bound is the classical Shannon source-coding inequality built
on Kraft's inequality \cite{Shannon1948,Kraft1949}. The proof is included here
explicitly because later formal coverage must match the actual rate argument,
not a citation boundary.
\end{remark}

\begin{theorem}[Finite seed charging]
\label{thm:finite-seed-charging}\lean{CoherenceCalculus.FiniteSeedChargingData, CoherenceCalculus.finite_seed_charging}
Suppose a seeded coded stochastic closure datum \(\mathbf C\) has finite seed
set
\[
\Xi=\{\xi_1,\dots,\xi_M\}.
\]
Then there exists a seeded coded stochastic closure datum
\(\overline{\mathbf C}\) with one-point seed space, alphabet
\[
\Xi\times\mathcal A
\]
and lengths
\[
\bar\ell(\xi_i,a):=\ell(a)+\lceil \log_2 M\rceil
\]
whose decoded output law is identical to that of \(\mathbf C\).
\end{theorem}

\begin{proof}
Let the new seed space be the singleton \(\{\ast\}\). Define the new encoder by
\[
\overline E_{x,\ast}(\xi_i,a):=\pi(\{\xi_i\})\,E_{x,\xi_i}(a),
\]
and define the new decoder by
\[
\overline D((\xi_i,a),\ast):=D(a,\xi_i).
\]
Then, for every Borel set \(B\subseteq V\),
\[
\kappa^{\overline{\mathbf C}}_x(B)
=
\sum_{i=1}^M\sum_{a\in\mathcal A}
\pi(\{\xi_i\})\,E_{x,\xi_i}(a)\,
\mathbf 1_B\!\bigl(D(a,\xi_i)\bigr)
=
\kappa^{\mathbf C}_x(B),
\]
so the decoded output law is unchanged. Moreover,
\[
\sum_{i=1}^{M}\sum_{a\in\mathcal A}2^{-\bar\ell(\xi_i,a)}
=
\frac{M}{2^{\lceil \log_2 M\rceil}}
\sum_{a\in\mathcal A}2^{-\ell(a)}
\le 1,
\]
so the new lengths satisfy Kraft's inequality.
\end{proof}

\begin{theorem}[Linear output composition]
\label{thm:linear-output-composition}\lean{CoherenceCalculus.LinearOutputCompositionData, CoherenceCalculus.linear_output_composition}
Let \(L:V\to W\) be linear, and let \(\mathbf C\) be a seeded coded stochastic
closure datum for a target \(T:K\to V\). Define \(L\mathbf C\) by keeping the
same seed, alphabet, lengths, and encoder, and replacing the decoder by
\[
D_{L\mathbf C}:=L\circ D.
\]
Then
\[
M_{L\mathbf C}=L\circ M_{\mathbf C},
\]
\[
\Gamma_{L\mathbf C}(x)=L\,\Gamma_{\mathbf C}(x)\,L^\ast,
\]
and the rate profiles are unchanged:
\[
L_K(L\mathbf C)=L_K(\mathbf C),
\qquad
L_{\mu}(L\mathbf C)=L_{\mu}(\mathbf C).
\]
\end{theorem}

\begin{proof}
All claims follow from linearity of expectation, the identity
\[
(Lu)\otimes(Lu)=L(u\otimes u)L^\ast,
\]
and the fact that encoder, alphabet, and code lengths are unchanged.
\end{proof}

\begin{theorem}[Unitary source/output covariance of coded closure]
\label{thm:unitary-coded-closure-covariance}\lean{CoherenceCalculus.UnitaryCodedClosureCovarianceData, CoherenceCalculus.unitary_coded_closure_covariance}
Let \(U:S\to S\) and \(W:V\to V\) be unitary, let \(K\subseteq S\) be a support
workload, and define
\[
UK:=\{Ux:x\in K\},
\qquad
T^{U,W}:UK\to V,
\qquad
T^{U,W}(x'):=WT(U^{-1}x')).
\]
Given a seeded coded stochastic closure datum \(\mathbf C\) for \(T\), define
\(\mathbf C^{U,W}\) on \(UK\) by
\[
\bigl(E^{U,W}\bigr)_{x',\xi}(a):=E_{U^{-1}x',\xi}(a),
\qquad
D^{U,W}(a,\xi):=W D(a,\xi).
\]
Then
\[
M_{\mathbf C^{U,W}}(Ux)=W M_{\mathbf C}(x),
\]
\[
\Gamma_{\mathbf C^{U,W}}(Ux)=W\Gamma_{\mathbf C}(x)W^\ast,
\]
and the pointwise rate profiles are preserved:
\[
L_{\mathbf C^{U,W}}(Ux)=L_{\mathbf C}(x).
\]
If \(\mu\) is a probabilistic workload datum with \(\operatorname{supp}\mu
\subseteq K\), then for the pushforward measure \(U_{\ast}\mu\) on \(UK\),
\[
L_{U_{\ast}\mu}(\mathbf C^{U,W})=L_{\mu}(\mathbf C),
\]
\[
\|M_{\mathbf C^{U,W}}-T^{U,W}\|_{L^2(U_{\ast}\mu;V)}
=
\|M_{\mathbf C}-T\|_{L^2(\mu;V)},
\]
and
\[
\int_{UK}\Tr\Gamma_{\mathbf C^{U,W}}(x')\,d(U_{\ast}\mu)(x')
=
\int_{K}\Tr\Gamma_{\mathbf C}(x)\,d\mu(x).
\]
\end{theorem}

\begin{proof}
For \(x'\in UK\), write \(x:=U^{-1}x'\). By definition,
\[
\widehat T_{\mathbf C^{U,W}}(x',\xi,a)
=
W\widehat T_{\mathbf C}(x,\xi,a).
\]
Linearity of expectation gives the mean identity, and the covariance identity
follows from the same tensor identity used in
Theorem~\ref{thm:linear-output-composition}. The pointwise rate profile is
unchanged because the transformed encoder uses the same symbol law and the same
length function at the corresponding preimage point.

If \(\mu\) is supported in \(K\), the rate identity for \(U_{\ast}\mu\) follows
from the pointwise equality and the change of variables \(x'=Ux\). The same
change of variables, together with the unitarity of \(W\), yields the
\(L^2\)-identity for the mean defect and the integrated covariance identity.
\end{proof}

\begin{definition}[Closure covariance operator]
\label{def:closure-covariance-operator}\lean{CoherenceCalculus.ClosureCovarianceOperatorData, CoherenceCalculus.closureCovarianceOperator}
Let \(\mathbf C\) be a seeded coded stochastic closure datum from a support
workload \(K\subseteq S\) to \(V\), and let \(\mu\) be a probabilistic
workload datum with \(\operatorname{supp}\mu\subseteq K\). The \emph{closure
covariance operator} is
\[
\mathfrak C_{\mathbf C,\mu}(T)
:=
\int_S \Gamma_{\mathbf C}(x)\,d\mu(x).
\]
\end{definition}

\begin{proposition}[Positivity of closure covariance]
\label{prop:closure-covariance-positive}\lean{CoherenceCalculus.ClosureCovarianceOperatorData, CoherenceCalculus.closure_covariance_positive}
In the setting of Definition~\ref{def:closure-covariance-operator}, the
operator \(\mathfrak C_{\mathbf C,\mu}(T)\) is positive semidefinite.
\end{proposition}

\begin{proof}
For every \(w\in V\),
\[
\langle w,\mathfrak C_{\mathbf C,\mu}(T)w\rangle
=
\int_S \langle w,\Gamma_{\mathbf C}(x)w\rangle\,d\mu(x)\ge 0,
\]
because each covariance operator \(\Gamma_{\mathbf C}(x)\) is positive
semidefinite.
\end{proof}

\begin{corollary}[Integrated stochastic closure error]
\label{cor:integrated-stochastic-closure-error}\lean{CoherenceCalculus.ClosureCovarianceOperatorData, CoherenceCalculus.integrated_stochastic_closure_error}
In the setting of Definition~\ref{def:closure-covariance-operator},
\[
\int_S \mathbb E\|\widehat T_{\mathbf C}(x,\xi,a)-T(x)\|^2\,d\mu(x)
=
\|M_{\mathbf C}-T\|_{L^2(\mu;V)}^2
+
\Tr\mathfrak C_{\mathbf C,\mu}(T).
\]
If \(\mathbf C\) is unbiased for \(T\), then
\[
\int_S \mathbb E\|\widehat T_{\mathbf C}(x,\xi,a)-T(x)\|^2\,d\mu(x)
=
\Tr\mathfrak C_{\mathbf C,\mu}(T).
\]
\end{corollary}

\begin{proof}
Apply Corollary~\ref{cor:integrated-bias-variance} to the induced kernel
\(\kappa^{\mathbf C}\). The second identity is the unbiased case.
\end{proof}

\begin{definition}[Observable coded stochastic returned closure]
\label{def:observable-coded-stochastic-returned-closure}\lean{CoherenceCalculus.ObservableReturnedClosureCode}
Let \(\Omega=(S,V,O)\) be an observable package on \(\cH\), let \(P\) be an
orthogonal projection, and fix \(z\). Set
\[
T_{\Omega,P,z}:=
\widetilde{\mathfrak R}_{\Omega}(P,A;z):S\to V.
\]
An \emph{observable coded stochastic returned closure datum} is a seeded coded
stochastic closure datum for the target transfer \(T_{\Omega,P,z}\). Its
corrected observable estimator is
\[
\widehat{\mathcal O}_{\mathbf C,\Omega,P,z}(x,\xi,a)
:=
O(PAPx)+z\,\widehat T_{\mathbf C}(x,\xi,a).
\]
\end{definition}

\begin{theorem}[Unbiased returned correction]
\label{thm:unbiased-returned-correction}\lean{CoherenceCalculus.ObservableReturnedClosureCode, CoherenceCalculus.unbiased_returned_correction}
In the setting of
Definition~\ref{def:observable-coded-stochastic-returned-closure}, if
\(\mathbf C\) is unbiased for \(T_{\Omega,P,z}\), then
\[
\mathbb E\bigl[
\widehat{\mathcal O}_{\mathbf C,\Omega,P,z}(x,\xi,a)
\bigr]
=
O(PAPx)+z\,\widetilde{\mathfrak R}_{\Omega}(P,A;z)x.
\]
Hence the exact Schur-corrected observable law is recovered in expectation.
\end{theorem}

\begin{proof}
By linearity of expectation and unbiasedness,
\[
\mathbb E\bigl[
\widehat{\mathcal O}_{\mathbf C,\Omega,P,z}(x,\xi,a)
\bigr]
=
O(PAPx)+z\,\mathbb E[\widehat T_{\mathbf C}(x,\xi,a)]
=
O(PAPx)+z\,T_{\Omega,P,z}x.
\]
\end{proof}

\begin{corollary}[Observable mean-square error]
\label{cor:observable-mean-square-error}\lean{CoherenceCalculus.ObservableReturnedClosureCode, CoherenceCalculus.observable_mean_square_error}
In the setting of Theorem~\ref{thm:unbiased-returned-correction}, for every
\(x\in S\),
\[
\mathbb E\Bigl[
\bigl\|
\widehat{\mathcal O}_{\mathbf C,\Omega,P,z}(x,\xi,a)
-
\bigl(O(PAPx)+z\,T_{\Omega,P,z}x\bigr)
\bigr\|^2
\Bigr]
=
|z|^2\,\mathbb E\|\widehat T_{\mathbf C}(x,\xi,a)-T_{\Omega,P,z}x\|^2.
\]
\end{corollary}

\begin{proof}
The deterministic term \(O(PAPx)\) cancels, and multiplication by the scalar
\(z\) scales the squared norm by \(|z|^2\).
\end{proof}

\begin{remark}[Boundary of the coded layer]
\label{rem:coded-closure-boundary}\lean{CoherenceCalculus.SeededClosureCode, CoherenceCalculus.ClosureWidthProfile, CoherenceCalculus.ObservableReturnedClosureCode}
The present stochastic/coded layer is still calculus-internal: it packages
kernel realizations, prefix-code lengths, entropy lower bounds, seed
accounting, and mean/covariance control for a fixed task-visible transfer.
Specific one-bit sketches, quantizers, sparse codes, or hardware schemes are
not new primitives here; they enter only as explicit admissible classes
\(\mathcal C\) inside this layer.
\end{remark}
